\section{Contexto mundial de la producci\'on y exportaci\'on de palta Hass}

\subsection{Importancia econ\'omica de la palta Hass}

La palta Hass (\textit{Persea americana} Mill.) se ha convertido en uno de los rubros agr\'icolas de mayor valor en el comercio internacional de frutas frescas, impulsada por su perfil nutricional, versatilidad gastron\'omica y creciente demanda en mercados de Norteam\'erica, Europa y Asia \citep{ADEX2023, VitteryFlores2023}. A diferencia de otros cultivos tradicionales, la palta combina un alto valor unitario con ventanas comerciales diferenciadas seg\'un origen geogr\'afico, lo que la convierte en un activo estrat\'egico para econom\'ias agr\'icolas exportadoras. \citet{SierraSelva2020} documenta que la producci\'on mundial super\'o las 7 millones de toneladas m\'etricas hacia finales de la d\'ecada de 2010, con una concentraci\'on creciente en variedades de exportaci\'on, especialmente Hass.

El peso econ\'omico del cultivo se refleja no solo en vol\'umenes, sino en el valor agregado de la cadena: empaque, certificaciones, log\'istica refrigerada y acceso a mercados premium. \citet{FAO2026} se\~nala que la palta encabeza el crecimiento de las frutas tropicales en t\'erminos de valor de exportaci\'on, superando a productos hist\'oricamente dominantes en ciertos segmentos. Para productores y exportadores, la palta Hass representa una fuente de divisas, empleo rural calificado e inversi\'on en infraestructura de riego tecnificado, aunque tambi\'en expone a los actores locales a din\'amicas globales de precios dif\'iciles de anticipar sin herramientas anal\'iticas.

La Tabla~\ref{tab:top_productores} resume la posici\'on relativa de los principales pa\'ises productores. M\'exico lidera ampliamente el ranking, seguido por Rep\'ublica Dominicana y Per\'u, lo que evidencia una geografia productiva diversa pero con concentraci\'on en Am\'erica Latina y algunos pa\'ises africanos y asi\'aticos. Esta distribuci\'on condiciona la competencia exportadora y la formaci\'on de precios internacionales.

\begin{table}[H]
\centering
\caption{Top 20 principales pa\'ises productores de palta --- Volumen (miles de toneladas m\'etricas)}
\label{tab:top_productores}
\small
\begin{tabular}{clr}
\toprule
\textbf{Rango} & \textbf{Pa\'is} & \textbf{Volumen (miles TM)} \\
\midrule
1  & M\'exico       & 2.184 \\
2  & Rep. Dominicana & 644 \\
3  & Per\'u         & 540 \\
4  & Colombia       & 324 \\
5  & Indonesia      & 304 \\
6  & Kenia          & 239 \\
7  & Chile          & 195 \\
8  & Brasil         & 183 \\
9  & Guatemala      & 168 \\
10 & Venezuela      & 152 \\
\bottomrule
\end{tabular}
\parbox{0.94\textwidth}{\small \textit{Nota.} Elaboraci\'on propia con base en \citep{SierraSelva2020}, Cuadro No.~1.}
\end{table}


\subsection{Crecimiento del consumo mundial}

El consumo global de palta ha experimentado un crecimiento sostenido vinculado con tendencias de alimentaci\'on saludable, mayor poder adquisitivo en mercados emergentes y popularizaci\'on del fruto en cadenas de restaurantes y retail \citep{ADEX2023, Zhang2025}. \citet{FAO2026} reporta incrementos significativos en vol\'umenes de importaci\'on en Estados Unidos, la Uni\'on Europea y mercados asi\'aticos, impulsados por campa\~nas de promoci\'on y mayor disponibilidad durante todo el a\~no gracias a la complementariedad estacional entre or\'igenes.

La Figura~\ref{fig:frutas_tropicales} sit\'ua a la palta dentro del panorama de frutas tropicales exportadas, evidenciando su dinamismo relativo frente a otros rubros. Paralelamente, la Figura~\ref{fig:proporciones_export_2025} muestra la composici\'on del comercio internacional por tipo y unidad de medida, lo cual resulta relevante para comprender c\'omo el valor y el volumen no evolucionan necesariamente en la misma proporci\'on.

\incluirfigura{figuras/figura1\_4.jpg}{0.88\linewidth}{Principales frutas tropicales: vol\'umenes totales de exportaci\'on agregados (2016--2025)}{fig:frutas_tropicales}

\incluirfigura{figuras/figura1\_5.jpg}{0.88\linewidth}{Proporciones del volumen de exportaci\'on por tipo (USD y toneladas), 2025}{fig:proporciones_export_2025}

La Tabla~\ref{tab:importaciones_pais} confirma que los principales mercados importadores concentran la demanda en Norteam\'erica y Europa, con participaci\'on creciente de Asia. Este patr\'on explica por qu\'e los productores peruanos deben monitorear no solo la oferta local, sino las se\~nales de absorci\'on en destinos clave y la competencia de otros or\'igenes en las mismas ventanas comerciales.

\begin{table}[H]
\centering
\caption{Evoluci\'on de las importaciones de palta por pa\'is (volumen en miles TM)}
\label{tab:importaciones_pais}
\small
\begin{tabular}{lrrrrr}
\toprule
\textbf{Pa\'is importador} & \textbf{2015} & \textbf{2016} & \textbf{2017} & \textbf{2018} & \textbf{2019} \\
\midrule
Estados Unidos & 890 & 920 & 980 & 1.050 & 1.120 \\
Pa\'ises Bajos  & 420 & 445 & 480 & 510 & 545 \\
Francia         & 180 & 195 & 210 & 225 & 240 \\
Reino Unido     & 165 & 172 & 185 & 198 & 215 \\
Jap\'on          & 120 & 128 & 135 & 142 & 155 \\
\bottomrule
\end{tabular}
\parbox{0.94\textwidth}{\small \textit{Nota.} Elaboraci\'on propia con base en \citep{SierraSelva2020}, Cuadro No.~13.}
\end{table}


\subsection{Principales pa\'ises productores y exportadores}

M\'exico mantiene el liderazgo productivo y exportador, apoyado en una base instalada consolidada, proximidad al mercado estadounidense y diversificaci\'on de destinos \citep{SierraSelva2020}. Per\'u, Chile, Colombia y Kenia han ampliado hect\'areas y vol\'umenes de forma acelerada, incrementando la presi\'on competitiva global. \citet{Trincado2022} analiza la complementariedad estacional entre Chile y Per\'u, subrayando que la superposici\'on parcial de campa\~nas puede generar picos de oferta y presi\'on a la baja sobre precios en determinados meses.

La Figura~\ref{fig:produccion_mundial} ilustra la proporci\'on de producci\'on mundial por pa\'is, visualizando la concentraci\'on latinoamericana. En exportaciones, la Tabla~\ref{tab:exportaciones_pais} muestra la evoluci\'on del valor en millones de d\'olares, donde Per\'u registra el crecimiento m\'as acelerado en el periodo 2015--2019, pasando de participaciones marginales a posicionarse entre los tres primeros exportadores globales.

\incluirfigura{figuras/figura1\_2.jpg}{0.85\linewidth}{Producci\'on mundial de palta en proporci\'on por pa\'is}{fig:produccion_mundial}

\begin{table}[H]
\centering
\caption{Evoluci\'on de las exportaciones de palta por pa\'is --- Valor (millones US\$)}
\label{tab:exportaciones_pais}
\small
\begin{tabular}{lrrrrr}
\toprule
\textbf{Pa\'is} & \textbf{2015} & \textbf{2016} & \textbf{2017} & \textbf{2018} & \textbf{2019} \\
\midrule
M\'exico   & 2.450 & 2.680 & 2.890 & 2.720 & 2.950 \\
Per\'u     &   580 &   720 &   890 & 1.120 & 1.380 \\
Chile      &   310 &   340 &   380 &   420 &   460 \\
Colombia   &   180 &   210 &   245 &   290 &   320 \\
Kenia      &   120 &   145 &   170 &   195 &   230 \\
\bottomrule
\end{tabular}
\parbox{0.94\textwidth}{\small \textit{Nota.} Elaboraci\'on propia con base en \citep{SierraSelva2020}, Cuadro No.~5.}
\end{table}


\subsection{Posicionamiento del Per\'u en el mercado internacional}

El Per\'u se ha consolidado como exportador estrat\'egico de palta Hass gracias a ventajas agroclim\'aticas en la costa norte y central, inversi\'on privada en agroexportaci\'on y articulaci\'on con operadores log\'isticos especializados \citep{VitteryFlores2023, Pantaleon2024}. \citet{VitteryFlores2023} sostienen que el pa\'is desplaz\'o progresivamente a competidores tradicionales y logr\'o ubicarse entre los principales proveedores globales, con especial protagonismo de la variedad Hass por su resistencia poscosecha y aceptaci\'on comercial.

Las exportaciones peruanas mantuvieron tendencia expansiva entre 2021 y 2025, alcanzando vol\'umenes superiores a 750 mil toneladas y valores FOB por encima de 1.300 millones USD en el \'ultimo a\~no reportado \citep{PromPeru2026, MIDAGRI2024}. No obstante, \citet{Capunay2025} advierten que el crecimiento exportador no garantiza rentabilidad uniforme para todos los eslabones de la cadena, especialmente cuando los precios unitarios caen por mayor competencia internacional.

La ventana comercial peruana, concentrada entre marzo y agosto, permite abastecer mercados cuando otros or\'igenes reducen disponibilidad, pero tambi\'en expone al sector a riesgos de sobreoferta si las decisiones de siembra no consideran proyecciones de precio. En este contexto, el posicionamiento del Per\'u como potencia exportadora refuerza la necesidad de herramientas predictivas que traduzcan datos de mercado en se\~nales comprensibles para productores.
